% ════════════════════════════════════════════════════════════
\section{CAPÍTULO III: METODOLOGÍA DE INVESTIGACIÓN}
% ════════════════════════════════════════════════════════════

\subsection{Técnicas e Instrumentos}
\subsubsection{Técnicas}
Minería de datos y análisis exploratorio de grafos (Data Mining \& Network Analysis), así como el modelado matemático computacional mediante cadenas de Markov en tiempo discreto.
\subsubsection{Instrumentos}
Lenguaje de programación Python 3.11, utilizando bibliotecas de análisis de redes computacionales (\texttt{NetworkX}, \texttt{python-igraph}) y procesamiento de álgebra lineal dispersa (\texttt{SciPy}, \texttt{NumPy}). Además, se utilizará el entorno de Jupyter Notebooks para la ejecución y visualización de resultados.

\subsection{Campo de Verificación}
\subsubsection{Ámbito Geográfico}
Debido a la naturaleza digital del problema, el ámbito no es geográfico físico, sino un ecosistema global de código abierto, específicamente el \textit{Python Package Index} (PyPI).
\subsubsection{Temporalidad}
Snapshot estático del ecosistema de dependencias correspondiente a la fecha de extracción del dataset público alojado en \textit{libraries.io}.

\subsection{Población y Muestra}
\subsubsection{Población}
La totalidad de los proyectos y paquetes de software registrados históricamente en PyPI (aproximadamente 500,000 paquetes).
\subsubsection{Muestra}
Subconjunto de paquetes activos seleccionados para el estudio. Se filtrarán aquellos paquetes con más de 100 descargas mensuales para garantizar la relevancia en entornos de producción, constituyendo una muestra de aproximadamente 100,000 paquetes activos.

\subsection{Recursos Necesarios}
\subsubsection{Recursos Humanos}
Los autores de la presente investigación.
\subsubsection{Recursos Materiales}
Estación de trabajo o computadora personal con entorno de ejecución de Python y suficiente capacidad de memoria RAM para almacenar el grafo disperso en memoria y procesar las iteraciones del algoritmo.
\subsubsection{Recursos Financieros}
Recursos propios (costos de energía eléctrica y acceso a internet). Herramientas y bibliotecas de código abierto utilizadas sin costo de licenciamiento.

\subsection{Estrategia de Recolección de Datos}

La información para esta investigación no se obtendrá mediante instrumentos cualitativos o encuestas de campo, sino a través de la extracción, procesamiento y estructuración de bases de datos públicas disponibles para la comunidad académica.

\subsubsection{Fuente de Datos: libraries.io – PyPI dependency graph}
Para este trabajo se utilizó la versión 1.6.0 del dataset público \textit{Libraries.io Open Source Repository and Dependency Metadata} \autocite{librariesio}, correspondiente a un snapshot publicado el 12 de enero de 2020. El mismo se encuentra disponible en la plataforma Zenodo (DOI: \texttt{10.5281/zenodo.3626071}) y se distribuye bajo la licencia \textit{Creative Commons Attribution-ShareAlike 4.0 International}.

A nivel global, este repositorio contiene información de 34 manejadores de paquetes distintos, distribuido en múltiples archivos CSV que superan los 120 GB en su tamaño original descomprimido. En particular, los datos maestros utilizados se concentran en \texttt{projects} (4,609,663 registros, 1.4 GB) y \texttt{dependencies} (190,388,563 dependencias declaradas, 21 GB).

Para el ecosistema de Python (PyPI), se extrajeron los siguientes campos relevantes:
\begin{itemize}
    \item \texttt{Project Name}: nombre del paquete dependiente.
    \item \texttt{Dependency Name}: nombre del paquete requerido.
    \item \texttt{Dependency Kind}: tipo de dependencia (\texttt{runtime},
          \texttt{development}, \texttt{optional}).
\end{itemize}
El análisis se restringe a dependencias de tipo \texttt{runtime}, que son las únicas relevantes para la propagación de fallos en producción. Luego del filtrado y la selección específica para el ecosistema de PyPI, el subconjunto final procesado en este estudio consta de 849,606 paquetes activos (nodos) y 6,976,793 dependencias de tipo \textit{runtime} (aristas).

\subsubsection{Pipeline de análisis}
\begin{enumerate}
    \item \textbf{Carga y preprocesamiento}: lectura del CSV de libraries.io;
          filtrado de dependencias \texttt{runtime}; eliminación de paquetes
          con menos de 100 descargas mensuales; construcción del grafo dirigido
          con \texttt{NetworkX}.
    \item \textbf{Análisis exploratorio topológico}: cálculo de distribución
          de grado (in/out), ajuste de ley de potencias, identificación de
          componentes fuertemente conexas, coeficiente de agrupamiento global.
    \item \textbf{Construcción de la cadena de Markov}: extracción de la
          matriz $\mathbf{H}$ en formato CSR (\textit{Compressed Sparse Row});
          identificación de nodos sumidero; construcción de $\mathbf{A}$ y
          $\mathbf{G}(d)$.
    \item \textbf{Cálculo del PageRank}: método de la potencia iterativo para
          $d \in \{0.75, 0.85, 0.90, 0.95\}$ hasta
          $\|\mathbf{r}^{(k+1)}-\mathbf{r}^{(k)}\|_1 < 10^{-8}$;
          registro de la curva de convergencia.
    \item \textbf{Detección de comunidades}: algoritmo de Louvain sobre el
          grafo no dirigido subyacente; asignación de cluster a cada paquete;
          visualización con layout ForceAtlas2.
    \item \textbf{Análisis comparativo}: correlación de Spearman $\rho_s$
          entre ranking PageRank y rankings por in-degree, descargas, y
          número de dependientes directos; identificación de paquetes con
          mayor divergencia entre PageRank e in-degree.
    \item \textbf{Análisis de cobertura transitiva}: para los Top-$k$
          paquetes por PageRank, cálculo del conjunto de paquetes que
          dependen de ellos transitivamente (BFS inverso), y estimación
          del porcentaje del ecosistema cubierto.
\end{enumerate}

\subsubsection{Implementación del Pipeline de Datos}
Para procesar el volumen masivo de datos (millones de nodos y aristas) sin comprometer la memoria, la implementación final del modelo no utilizó scripts monolíticos, sino un \textit{pipeline} estructurado en etapas:

\begin{itemize}
    \item \textbf{Procesamiento de Datos (ETL):} Los datos brutos originales fueron convertidos al formato columnar Apache Parquet utilizando la librería \texttt{pyarrow}. Esto permitió cargar y filtrar eficientemente los millones de registros (filtrando exclusivamente el ecosistema PyPI y las dependencias de ejecución real o \textit{runtime}) en tiempos de milisegundos usando \texttt{pandas}.
    \item \textbf{Construcción del Grafo y Matrices:} Se utilizó \texttt{networkx} para modelar la topología de la red dirigida. Para realizar los cálculos algebraicos, la matriz de adyacencia se tradujo a una matriz rala (Sparse Matrix) en formato CSR (Compressed Sparse Row) mediante \texttt{scipy.sparse}, lo que redujo drásticamente el consumo de memoria RAM.
    \item \textbf{Método de la Potencia Vectorizado:} La iteración de la Matriz de Google no se realizó con bucles \texttt{for} estándar, sino a través de operaciones matriciales vectorizadas utilizando \texttt{numpy}. El manejo explícito de los "nodos sumidero" (paquetes sin dependencias) se resolvió algebraicamente sin modificar la matriz rala original, preservando la eficiencia temporal.
    \item \textbf{Detección de Comunidades:} Para identificar los clusters semánticos, se integró la librería \texttt{python-louvain} (basada en heurísticas de optimización de modularidad), ejecutada sobre la proyección no dirigida del grafo de dependencias.
\end{itemize}

Todo el código fuente desarrollado para este trabajo, incluyendo los scripts de preprocesamiento, análisis de grafos y generación de visualizaciones, se encuentra estructurado de manera modular y es completamente reproducible.
